\section{Lựa chọn mô hình triển khai}

Model triển khai được chọn theo hai bước: loại các cấu hình vi phạm ràng buộc
độ chính xác, latency hoặc ổn định nhiệt; trong nhóm còn lại, ưu tiên cấu hình có
latency thấp nhất trên thiết bị yếu hơn (Jetson Nano) vì đó là giới hạn của cả đội
thiết bị.

Dựa trên phân tích độ phức tạp và kết quả thực nghiệm sơ bộ (Mục~\ref{sec:pilot-results}),
cấu hình được đề xuất là \textbf{YOLO26n cắt tỉa 30\% + phục hồi + QAT, chạy
TensorRT INT8} (Bảng~\ref{tab:selected-model}). Lý do: cắt tỉa giảm 21\% FLOPs,
INT8 giảm tiếp chi phí mỗi phép toán và bộ nhớ trọng số khoảng 4 lần so với
FP32, còn QAT giữ độ chính xác tốt hơn PTQ. Với thiết bị hỗ trợ kém INT8, engine
FP16 của mô hình đã cắt tỉa là phương án dự phòng. Kết luận cuối cùng phụ thuộc
số đo trên thiết bị ở Chương~6.

\begin{table}[H]
\centering
\caption{Cấu hình model đề xuất triển khai}
\label{tab:selected-model}
\begin{tabularx}{\textwidth}{L{4cm}Y}
\toprule
\textbf{Thuộc tính} & \textbf{Giá trị}\\
\midrule
Kiến trúc & YOLO26n, đầu vào $640\times640$, 11 lớp PPE\\
Phương pháp tối ưu & SHO baseline $\rightarrow$ sparsity $\rightarrow$ pruning 30\% $\rightarrow$ chưng cất $\rightarrow$ QAT\\
Độ phức tạp & 2,17 M tham số, 4,82 GFLOPs\\
Precision & INT8 (TensorRT), dự phòng FP16\\
Định dạng triển khai & TensorRT \texttt{.engine} build riêng cho từng loại Jetson\\
Phân phối & Kho model Hugging Face Hub, ghim theo commit SHA\\
Độ chính xác & \cbs{mAP@50, mAP@50:95 trên tập test}\\
Hiệu năng trên thiết bị & \cbs{FPS, latency trên Jetson Nano và Orin Nano}\\
\bottomrule
\end{tabularx}
\end{table}
